% --- The broader world-model landscape ---
\begin{frame}{Zooming Out: The World-Model Landscape}
  We have followed two lineages in depth (RSSM/Dreamer and JEPA). The wider field sorts by two questions: \textbf{what does the model predict}, and \textbf{what is it for}?
  \begin{itemize}\itemsep0.24em
    \item \textbf{Latent dynamics for control:} PlaNet, Dreamer v1--v3, TD-MPC2. Predict compact latents plus reward; train policies in imagination.
    \item \textbf{Value-equivalent planning:} MuZero. Predict only what planning needs (reward, value, policy), nothing about appearance.
    \item \textbf{Latent prediction (JEPA):} V-JEPA 2-AC, DINO-WM, PLDM, LeWM. Predict representations; plan by latent search.
    \item \textbf{Video simulators:} Sora, GAIA-1/2, NVIDIA Cosmos. Generate future video, often for data and evaluation.
    \item \textbf{Interactive generative worlds:} Genie 1--3, GameNGen, Oasis. Generate the next frame \emph{conditioned on user actions}, in real time.
    \item \textbf{Spatial intelligence / 3D worlds:} World Labs (Fei-Fei Li). Generate explicit, persistent 3D scenes rather than video streams.
  \end{itemize}
\end{frame}

\begin{frame}{Roots: Schmidhuber's Controller--Model Systems and Dyna}
  \begin{itemize}\itemsep0.28em
    \item \textbf{Schmidhuber (1990):} a recurrent \textbf{controller} network paired with a recurrent \textbf{model} network that learns to predict the consequences of actions; the controller can be trained through the differentiable learned model. Essentially the 2018 V--M--C blueprint, three decades earlier.
    \item \textbf{Artificial curiosity (1991):} reward the agent where the model's \emph{prediction error} is high, so it explores what it does not yet understand; the world model doubles as an exploration signal.
    \item \textbf{Sutton's Dyna (1991):} in parallel, RL work framed planning as \emph{learning from imagined experience} generated by a learned model, the same idea Dreamer scales up.
    \item Takeaway for the exam: ``training a policy inside a learned model'' is not a 2018 invention; 2018--2025 made it \textbf{work at scale} with deep networks.
  \end{itemize}
\end{frame}

\begin{frame}{MuZero: Planning with a Value-Equivalent Model}
  \begin{itemize}\itemsep0.28em
    \item \textbf{MuZero} (Schrittwieser et al., Nature 2020) learns a model that predicts only three things: \textbf{reward}, \textbf{value}, and \textbf{policy}. It never reconstructs observations at all.
    \item The model is \textbf{value-equivalent}: it only has to be right about quantities that matter for decisions, not about how the world looks.
    \item Planning: Monte Carlo tree search unrolled in the learned latent space; reached superhuman play in Go, chess, shogi, and strong Atari results without being given the rules.
    \item Contrast with today's other families: Dreamer keeps a decoder (reconstruction grounds the latent), JEPA drops the decoder but still predicts \emph{representations}; MuZero is the extreme end where the ``world'' is reduced to its decision-relevant summary.
  \end{itemize}
\end{frame}

\begin{frame}{Interactive Generative Worlds: GameNGen, Oasis, Genie}
  \begin{itemize}\itemsep0.26em
    \item \textbf{GameNGen} (Google, 2024): a diffusion model that runs \emph{DOOM} at interactive frame rates; the game engine is replaced entirely by next-frame prediction conditioned on player actions.
    \item \textbf{Oasis} (Decart and Etched, 2024): a playable Minecraft-style world generated frame by frame in real time from keyboard and mouse input.
    \item \textbf{Genie} (Google DeepMind): a foundation world model trained on internet gameplay video; it learns \textbf{latent actions} without any action labels, so generated worlds become playable. Genie 2 (2024) produced 3D playable scenes from one image; \textbf{Genie 3} (2025) generates interactive worlds from a text prompt in real time at 720p, staying consistent for minutes (as reported by DeepMind).
    \item These systems answer the world-model question $p(s' \mid s, a)$ directly \textbf{in pixel space}, with the user supplying $a$ live.
  \end{itemize}
\end{frame}

\begin{frame}{Embodied and Driving Simulators: GAIA and Cosmos}
  \begin{itemize}\itemsep0.28em
    \item \textbf{GAIA-1/GAIA-2} (Wayve): generative world models for driving; produce controllable, multi-camera driving scenarios (weather, agents, rare events) used to train and evaluate autonomous-driving stacks.
    \item \textbf{NVIDIA Cosmos} (2025): a \emph{world foundation model} platform for ``Physical AI''; pretrained video world models released with open weights, meant to be post-trained into robot- or vehicle-specific simulators for synthetic data generation and policy evaluation.
    \item The economics: real robot and driving data is expensive and dangerous to collect at the tails; a world model that generates rare, safety-critical scenarios earns its keep even if it is never used for planning directly.
    \item Same generative-simulator recipe as Sora, pointed at embodiment rather than media.
  \end{itemize}
\end{frame}

\begin{frame}{Spatial Intelligence: Fei-Fei Li and World Labs}
  \begin{itemize}\itemsep0.26em
    \item \textbf{Thesis (Fei-Fei Li):} language is not enough; \textbf{spatial intelligence}, i.e.\ perceiving, reasoning about, and generating 3D space, is the next frontier for AI. World Labs (founded 2024) builds toward it.
    \item \textbf{RTFM} (2025): a real-time frame model that renders explorable views of a scene interactively on modest hardware (as reported by World Labs).
    \item \textbf{Marble} (2025): generates \textbf{persistent, editable 3D worlds} (Gaussian-splat scenes) from text, images, or video; worlds can be exported into standard graphics and game pipelines.
    \item Key difference from video world models: the world is an \textbf{explicit 3D state}, so it does not drift or forget; geometry persists when you look away. The cost is that dynamics and agents must be added on top.
    \item For robotics and AR, this line argues the world model should \emph{be} a 3D scene, not a video stream.
  \end{itemize}
\end{frame}

\begin{frame}{The Landscape at a Glance}
  \footnotesize
  \begin{center}
    \renewcommand{\arraystretch}{1.3}
    \begin{tabular}{@{}p{2.8cm}p{3.1cm}p{3.6cm}p{3.0cm}@{}}
      \hline
      \textbf{Family} & \textbf{Predicts} & \textbf{Representatives} & \textbf{Primary use} \\
      \hline
      Latent dynamics (RSSM) & latents + reward (pixels optional; TD-MPC2 drops them) & PlaNet, Dreamer v1--v3, TD-MPC2 & policy learning in imagination \\
      Value-equivalent & reward, value, policy only & MuZero & search-based planning \\
      Latent prediction (JEPA) & representations & V-JEPA 2-AC, DINO-WM, PLDM, LeWM & latent MPC, robotics \\
      Video simulators & future video & Sora, GAIA-1/2, Cosmos & synthetic data, evaluation \\
      Interactive worlds & next frame given action & Genie 3, GameNGen, Oasis & playable environments \\
      3D scene generation & persistent 3D scenes & World Labs Marble, RTFM & AR, games, robotics \\
      \hline
    \end{tabular}
  \end{center}
  \vspace{0.3em}
  The families increasingly borrow from each other: interactive worlds add memory, latent planners add dense features, and simulators add action conditioning.
\end{frame}
